\chapter*{Annexe A : Catalogue de l'API RESTful}
\addcontentsline{toc}{chapter}{Annexe A : Catalogue de l'API RESTful}
\markboth{Annexe A : Catalogue de l'API RESTful}{Annexe A : Catalogue de l'API RESTful}

\par Cette premi\`ere annexe d\'etaille le catalogue exhaustif des points d'entr\'ee de l'API RESTful expos\'es par la plateforme, avec leurs param\`etres de filtrage et les formats de r\'eponse normalis\'es.

\begin{table}[H]
\centering
\caption{Catalogue Exhaustif des Endpoints de l'API RESTful}
\label{tab:annexe_api_catalog}
\resizebox{\textwidth}{!}{%
\begin{tabular}{|L|C|L|L|C|}
\hline
\textbf{Endpoint} & \textbf{Verbe} & \textbf{Param\`etres Cl\'es} & \textbf{Description} & \textbf{Code} \\
\hline
\texttt{/api/v1/auth/login} & POST & \texttt{email}, \texttt{password} & Authentification et g\'en\'eration du token Sanctum & 200 \\
\hline
\texttt{/api/v1/projects} & GET & \texttt{page}, \texttt{search}, \texttt{status} & Liste pagin\'ee des projets avec m\'etriques d'audit & 200 \\
\hline
\texttt{/api/v1/projects} & POST & \texttt{name}, \texttt{targets[]}, \texttt{poa\_file} & Cr\'eation de projet avec validation de mandat PoA & 201 \\
\hline
\texttt{/api/v1/scans} & POST & \texttt{project\_id}, \texttt{type}, \texttt{profile} & Mise en file d'attente asynchrone (Redis Streams) d'un scan & 202 \\
\hline
\texttt{/api/v1/scans/\{id\}} & GET & \texttt{id} & R\'ecup\'eration de l'\'etat et des r\'esultats du scan & 200 \\
\hline
\texttt{/api/v1/findings} & GET & \texttt{severity}, \texttt{project\_id}, \texttt{cve} & Liste filtr\'ee des constats normalis\'es & 200 \\
\hline
\texttt{/api/v1/findings/\{id\}/remediate} & POST & \texttt{type} (bash/ansible/docker) & G\'en\'eration de script Remediation-as-Code via LLM & 200 \\
\hline
\texttt{/api/v1/projects/\{id\}/graph} & GET & \texttt{depth}, \texttt{filter\_severity} & Donn\'ees s\'erialis\'ees du graphe Apache AGE & 200 \\
\hline
\texttt{/api/v1/assets/\{id\}/impact} & GET & \texttt{k\_max} (d\'efaut: 3) & Calcul algorithmique du rayon d'impact (BFS) & 200 \\
\hline
\texttt{/api/v1/reports/\{id\}/generate} & POST & \texttt{format} (pdf/json/zip) & Compilation asynchrone du rapport d'expertise & 200 \\
\hline
\texttt{/api/v1/audit-logs} & GET & \texttt{user\_id}, \texttt{action}, \texttt{date} & Pistes d'audit immuables pour contr\^ole conformit\'e & 200 \\
\hline
\end{tabular}}%
\end{table}

\clearpage

\chapter*{Annexe B : Art\'efacts DevSecOps}
\addcontentsline{toc}{chapter}{Annexe B : Art\'efacts DevSecOps}
\markboth{Annexe B : Art\'efacts DevSecOps}{Annexe B : Art\'efacts DevSecOps}

\par Cette annexe documente les art\'efacts d'infrastructure et d'int\'egration continue mis en \oe uvre pour garantir la s\'ecurit\'e intrins\`eque du syst\`eme.

\section*{B.1 Extrait de Configuration Docker Compose Durcie}
\par D\'eclaration du microservice de reconnaissance au sein de \texttt{docker-compose.prod.yml} :
{\footnotesize
\begin{verbatim}
services:
  recon-service:
    image: ghcr.io/aymenazizi/cybersec-recon:v1.0.0
    container_name: cybersec-recon
    restart: unless-stopped
    user: "1001:1001"
    read_only: true
    cap_drop:
      - ALL
    security_opt:
      - no-new-privileges:true
    tmpfs:
      - /tmp:size=128M,noexec,nosuid,nodev,mode=1777
    networks:
      - cybersec-internal
    environment:
      - FLASK_ENV=production
      - REDIS_HOST=redis
      - DB_HOST=postgres
    deploy:
      resources:
        limits:
          cpus: '1.5'
          memory: 1024M
\end{verbatim}
}

\section*{B.2 Workflow DevSecOps sous GitHub Actions}
\par Extrait du pipeline CI/CD automatis\'e (\texttt{.github/workflows/devsecops.yml}) :
{\footnotesize
\begin{verbatim}
name: DevSecOps Pipeline
on:
  push:
    branches: [ main ]
  pull_request:
    branches: [ main ]

jobs:
  sast-sca-sbom:
    runs-on: ubuntu-latest
    steps:
      - uses: actions/checkout@v4
      - name: PHP Static Analysis (Psalm)
        run: ./vendor/bin/psalm --config=psalm.xml --level=3
      - name: Python Static Analysis (Bandit)
        run: bandit -r ./services/ -ll -ii
      - name: Software Composition Analysis (Trivy FS)
        uses: aquasecurity/trivy-action@master
        with:
          scan-type: 'fs'
          severity: 'CRITICAL,HIGH'
          exit-code: '1'
      - name: Generate CycloneDX SBOM (Syft)
        uses: anchore/sbom-action@v0
        with:
          format: cyclonedx-json
          output-file: sbom.cyclonedx.json

  build-sign-deploy:
    needs: sast-sca-sbom
    runs-on: ubuntu-latest
    steps:
      - name: Build and Push OCI Image
        uses: docker/build-push-action@v5
        with:
          push: true
          tags: ghcr.io/aymenazizi/cybersec-platform:latest
      - name: Sign Container Image (Cosign)
        run: cosign sign --yes ghcr.io/aymenazizi/cybersec-platform:latest
\end{verbatim}
}

\section*{B.3 Sch\'ema Pydantic pour l'IA Contrainte}
{\footnotesize
\begin{verbatim}
class RemediationOutputSchema(BaseModel):
    finding_id: int = Field(..., description="ID unique du constat")
    remediation_type: Literal["bash", "ansible", "dockerfile", "terraform"]
    summary: str = Field(..., max_length=255, description="Titre")
    risk_mitigated: str = Field(..., description="Risque atténué")
    script_content: str = Field(..., description="Code exécutable")
    evidence: EvidenceCitationSchema = Field(..., description="Citations")
    safety_notes: Optional[str] = Field(None, description="Notes")
\end{verbatim}
}

\clearpage

\chapter*{Annexe C : Nginx et Requ\^etes openCypher}
\addcontentsline{toc}{chapter}{Annexe C : Nginx et Requ\^etes openCypher}
\markboth{Annexe C : Configuration Nginx et openCypher}{Annexe C : Configuration Nginx et openCypher}

\section*{C.1 Configuration Nginx Durcie avec En-t\^etes de S\'ecurit\'e}
{\footnotesize
\begin{verbatim}
server {
    listen 443 ssl http2;
    server_name aymenazizi.dijaly.com;

    ssl_certificate /etc/letsencrypt/live/aymenazizi.dijaly.com/fullchain.pem;
    ssl_certificate_key /etc/letsencrypt/live/aymenazizi.dijaly.com/privkey.pem;
    ssl_protocols TLSv1.3;
    ssl_ciphers HIGH:!aNULL:!MD5;

    # En-têtes de sécurité obligatoires
    add_header Strict-Transport-Security 
               "max-age=31536000; includeSubDomains" always;
    add_header X-Frame-Options "DENY" always;
    add_header X-Content-Type-Options "nosniff" always;
    add_header Referrer-Policy "strict-origin-when-cross-origin" always;
    add_header Content-Security-Policy 
               "default-src 'self'; script-src 'self' 'unsafe-inline';";

    location / {
        proxy_pass http://cybersec-backend:8000;
        proxy_set_header Host $host;
        proxy_set_header X-Real-IP $remote_addr;
        proxy_set_header X-Forwarded-For $proxy_add_x_forwarded_for;
        proxy_set_header X-Forwarded-Proto https;
    }
}
\end{verbatim}
}

\section*{C.2 Requ\^etes openCypher Avanc\'ees dans Apache AGE}
\par Extraction des chemins d'attaque reliant un actif racine \`a une vuln\'erabilit\'e critique :
{\footnotesize
\begin{verbatim}
SELECT * FROM cypher('cybersec_graph', $$
  MATCH path = (t:Target)-[:RESOLVES_TO]->(a:Asset)-[:HOSTS]->(p:Port)
               -[:RUNS]->(s:Service)-[:HAS_VULN]->(v:Vulnerability)
  WHERE v.severity = 'CRITICAL'
  RETURN t.name, a.name, p.number, s.name, v.cve_id, v.cvss_score
  ORDER BY v.cvss_score DESC
$$) as (target varchar, asset varchar, port int, 
        service varchar, cve varchar, cvss float);
\end{verbatim}
}

\section*{C.3 Playbook Ansible de Rem\'ediation SSH G\'en\'er\'e par l'IA}
{\footnotesize
\begin{verbatim}
---
- name: Hardening Configuration SSH Serveur
  hosts: all
  become: yes
  tasks:
    - name: Désactiver l'authentification par mot de passe
      lineinfile:
        path: /etc/ssh/sshd_config
        regexp: '^#?PasswordAuthentication'
        line: 'PasswordAuthentication no'
        state: present

    - name: Interdire la connexion directe du compte root
      lineinfile:
        path: /etc/ssh/sshd_config
        regexp: '^#?PermitRootLogin'
        line: 'PermitRootLogin prohibit-password'
        state: present

    - name: Redémarrer le démon SSH de façon sécurisée
      service:
        name: sshd
        state: restarted
\end{verbatim}
}
